% Slajd 2 – problem i grupa docelowa (kryterium: związek z wyzwaniem i użyteczność, 25 %).
\begin{frame}{„Dostępne / niedostępne” to za mało}
  \begin{columns}[T,onlytextwidth]
    \begin{column}{0.49\textwidth}
      \begin{block}{Problem}
        \begin{itemize}
          \item Jedna etykieta nie mówi, czy wjadę z~wózkiem, czy drzwi mają 90~cm, czy jest toaleta i~gdzie można odpocząć.
          \item Informacje w~sieci nie mają źródła ani daty: zepsuta winda zamienia „dostępne” w~zmarnowaną wyprawę.
          \item Brak informacji bywa pokazywany tak, jakby był jej potwierdzeniem.
          \item Miasto nie ma zasobów, by ręcznie utrzymywać bazę -- dane muszą aktualizować się same.
        \end{itemize}
      \end{block}
    \end{column}
    \begin{column}{0.49\textwidth}
      \begin{block}{Dla kogo -- zakres prototypu}
        \begin{itemize}
          \item \textbf{Grupa główna:} osoby poruszające się na wózkach oraz rodzice z~wózkami dziecięcymi (profile „Wózek” i~„Wózek dziecięcy”).
          \item Te same dane obsługują kolejne profile: Słabo widzę, Senior, Po zabiegu (turystyka medyczna), Bagaż, Podróżuję z~psem, Rower.
          \item Mieszkańcy i~turyści (interfejs PL, EN, DE, UK) oraz \textbf{właściciele i~zarządcy obiektów}, którzy dostają własny panel.
        \end{itemize}
      \end{block}
    \end{column}
  \end{columns}
  \vspace{2mm}
  \begin{exampleblock}{Nasza odpowiedź}
    Accessly odpowiada na pytanie \textbf{„czy to miejsce zadziała dla mnie?”} konkretnymi faktami o~barierach
    i~udogodnieniach. Każdy fakt ma \textbf{źródło, datę i~status weryfikacji}, a~brak danych jest pokazany
    jako brak danych -- nigdy jako dostępność.
  \end{exampleblock}

  \note[item]{Zacząć od historii: Anna na wózku chce wyjść na kawę w~centrum. Mapa z~zeszłego roku mówi „dostępne”, a~dziś przy wejściu jest próg i~nikt nie wie, czy toaleta jest na dole.}
  \note[item]{Podkreślić dwie grupy: użytkownicy końcowi oraz właściciele – bez nich dane nie będą aktualne.}
  \note[item]{Profil to preferencje (schody, winda, cisza), nie informacja o~niepełnosprawności – wraca na slajdzie o~prywatności.}
\end{frame}
